\section{Struttura delle simulazioni (NED, INI, SUMO)}

Una volta completata la compilazione, la configurazione sperimentale è stata organizzata lungo tre assi:
\begin{itemize}
\item topologia e moduli OMNeT++ (file NED);
\item parametrizzazione degli esperimenti (file \texttt{omnetpp.ini});
\item mobilità microscopica (file SUMO: \texttt{.sumocfg}, \texttt{.net.xml}, \texttt{.rou.xml}).
\end{itemize}

La separazione tra questi livelli consente di modificare scenari e politiche senza alterare il codice C++ dei moduli.

\begin{lstlisting}[caption={Schema semplificato di rete NED per co-simulazione},label={lst:ned_skeleton}]
network PlatooningScenario
{
    submodules:
        manager: TraCIScenarioManagerLaunchd;
        node[NUM_CARS]: Car;
        rsu: RoadSideUnit;
}
\end{lstlisting}

\begin{lstlisting}[caption={Estratto tipo di configurazione INI multi-interfaccia},label={lst:ini_skeleton}]
[General]
network = PlatooningScenario
sim-time-limit = 200s

# OMNeT++ / Veins / TraCI
*.manager.updateInterval = 0.1s
*.manager.host = "localhost"
*.manager.moduleType = "org.car2x.plexe.Car"

# Output
output-scalar-file = ${resultdir}/${configname}-${runnumber}.sca
output-vector-file = ${resultdir}/${configname}-${runnumber}.vec
\end{lstlisting}

\begin{lstlisting}[caption={Estratto tipo di configurazione SUMO},label={lst:sumo_cfg_skeleton}]
<configuration>
  <input>
    <net-file value="highway.net.xml"/>
    <route-files value="platoon.rou.xml"/>
  </input>
  <time>
    <step-length value="0.1"/>
  </time>
</configuration>
\end{lstlisting}

Nel workflow adottato, la coerenza tra INI e SUMO è stata verificata su tre punti: identificativi dei veicoli, durata della simulazione e passo temporale. In assenza di tale allineamento, la simulazione può avviarsi correttamente ma produrre traiettorie o metriche non confrontabili tra run.
